\section{Discussion and conclusion}
\label{sec:conclusion}

A counterfactual explanation is advice, and advice is only worth giving if it can
be followed and would work. The field has largely measured whether advice flips
the model. We argued that this conflates three separate things --- flipping the
model, surviving the causal consequences of the change, and being something a
person could actually do --- and that only their conjunction, \Usable{}, is what a
rejected applicant needs. Measuring the three separately exposes a gap that
validity-only evaluation cannot see: across four public datasets every method
clears validity, yet usable recourse ranges from $13\%$ to $100\%$, and the gap
opens precisely where the features are densely causal. \cac{} closes it by
treating causal validity as structural --- it searches for interventions and
scores them through a structural causal model, so it can only ever propose a world
its own model endorses --- and an ablation shows the causal term is specifically
what closes it, at a modest and honestly-reported cost in effort.

\paragraph{Limitations.} We are candid about three. \emph{The causal models are
hand-crafted and imperfect}: the diabetes graph is well supported by the data, but
the German and HELOC graphs leave many real dependences unmodelled
(\S\ref{sec:scm}), and \cac{}'s guarantees are relative to the graph it is given
--- mitigated, not erased, by the graceful degradation we measured under graph
perturbation (\S\ref{sec:exp:robust}), and by the model-transfer study showing the
advantage does not depend on scoring against the same box \cac{} searched.
\emph{The human-feasibility audit is a model, not people}: we encode mutability
and effort from domain knowledge and ship a user-study protocol (Appendix and
\texttt{user\_study/}) to test those assumptions with real respondents, but we do
not run it here. \emph{Recourse is single-shot}: we do not model the applicant
taking one action, being re-scored, and acting again.

\paragraph{Reproducibility.} Every dataset is public and fetched from Kaggle or
\texttt{ucimlrepo}; every model trains locally; no result depends on a paid
service or a network call at evaluation time. The pipeline is deterministic
(\texttt{random\_state}$=42$) and rebuilds end to end --- data, models, SCMs, cost
models, all methods, every table and figure in this paper, and the demo --- from a
single \texttt{make all}, with \texttt{make paper} producing this PDF.

\paragraph{Outlook.} The ingredients of good recourse --- validity, causal
consequence, actionability --- were already in the literature separately. The step
that mattered was refusing to trade them off: auditing all three at once, and
building a method whose search space cannot contain a causally impossible plan.
The result is recourse that survives its own consequences. We hope the three-audit
evaluation, more than the method, is what the field takes forward: it is cheap to
adopt, it applies to any recourse method, and it turns ``does it flip the box'' back
into the question that matters --- ``could this person follow your advice, and
would it work.''
